\subsection{Estructura de costes}

\begin{table}[H]
\centering
\small
\rowcolors{1}{}{ }
\begin{tabularx}{\textwidth}{
>{\columncolor{gray!15}}X
>{\columncolor{white}}c
>{\columncolor{gray!15}}c
>{\columncolor{white}}c
}
\hline
\textbf{Componentes} &
\textbf{Precio / unidad (\texteuro)} &
\textbf{Nº unidades} &
\textbf{Precio proyecto (\texteuro)} \\
\hline
Arduino Mega & 25,99 & 1 & 25,99 \\
Driver L298N & 1,68 & 2 & 3,36 \\
Motor DC & 2,34 & 4 & 9,36 \\
Sensor paso carta & 1,34 & 3 & 4,02 \\
Sensor magnético home & 0,88 & 1 & 0,88 \\
Pantalla LCD con pulsadores & 4,18 & 1 & 4,18 \\
Batería 9 V & 8,00 & 1 & 8,00 \\
Imanes & 0,12 & 6 & 0,72 \\
Rodillos cartas & 0,60 & 6 & 3,60 \\
Anillo colector (Slip Ring) & 4,00 & 1 & 4,00 \\
Rodamientos y ejes & -- & 3 & 1,00 \\
Tornillos & -- & -- & 0,00 \\
Tuercas & -- & -- & 0,00 \\
Varillas & -- & -- & 0,00 \\
Impresión 3D & -- & 1 & 0,00 \\
\hline
\multicolumn{3}{r}{\textbf{TOTAL}} & \textbf{61,51} \\
\hline
\end{tabularx}
\caption{Coste de los componentes del prototipo del proyecto}
\label{tab:costes}
\end{table}


\subsection{Producción}
El sistema desarrollado se ha concebido como un prototipo funcional basado en la integración de componentes comerciales estándar y elementos mecánicos personalizados. La producción se ha planteado mediante un ensamblaje modular, permitiendo el montaje y verificación independiente de los subsistemas electrónico, mecánico y de control.

La electrónica se basa en una placa Arduino Mega junto con sensores, actuadores y módulos de potencia seleccionados por su bajo coste, disponibilidad comercial y facilidad de integración. El diseño mecánico emplea piezas fabricadas mediante impresión 3D y elementos estándar, lo que facilita una fabricación rápida y la modificación del prototipo. En un escenario de producción en serie, estas piezas podrían sustituirse por procesos industriales para reducir costes y aumentar la robustez.

El proceso de producción no requiere maquinaria especializada ni personal altamente cualificado, lo que permite una fácil escalabilidad y refuerza la viabilidad técnica y económica del proyecto.

\subsection{Viabilidad económica}
El análisis de la viabilidad económica del proyecto se ha realizado a partir de la estimación detallada de los costes de los componentes necesarios para la implementación del prototipo, tal y como se refleja en la Tabla~\ref{tab:costes}. El coste total del conjunto de elementos hardware asciende a \textbf{61,51~\texteuro}, incluyendo la electrónica de control, sensores, actuadores y elementos mecánicos auxiliares. Algunos elementos, como la batería, los rodillos o la impresión 3D, no suponen un coste directo al tratarse de materiales reutilizados o fabricados internamente.

Además, en un escenario de producción a mayor escala, el coste por unidad podría verse reducido mediante compras al por mayor y la optimización del diseño, situándose el coste de fabricación en un rango estimado de \textbf{35--50~\texteuro} por unidad.

A partir de estos costes, sería viable establecer un \textit{Precio de Venta al Público} (PVP) competitivo, comprendido entre \textbf{89~\texteuro\ y 129~\texteuro}, lo que permitiría cubrir los gastos asociados a fabricación, distribución y comercialización, así como obtener un margen de beneficio razonable que haga sostenible el producto.

En conclusión, el proyecto presenta una \textbf{viabilidad económica favorable}, al combinar un coste de producción contenido con la posibilidad de fijar un precio competitivo en el mercado, lo que refuerza su potencial de explotación y comercialización.
